% =============================================================================
% CHAPTER 3: SYSTEM DESIGN AND ANALYSIS
% =============================================================================

\chapter{System Design and Analysis}

This chapter presents the design and implementation of the Dynamic Context-Aware Trie (DCA-Trie) framework. The chapter begins with an overview of the full system architecture in Section~\ref{sec:architecture}, situating each component within the end-to-end pipeline. Section~\ref{sec:oracle} details the constraint oracle layer, which is the primary contribution of this work, covering the formal problem definition, the semantic relevance scorer, and the Semantic Irrelevance Ratio metric. Section~\ref{sec:v1} describes DCA-Trie v1, the static semantic filtering variant. Section~\ref{sec:v2} describes DCA-Trie v2, the step-wise dynamic expansion variant. Section~\ref{sec:integration} explains how both variants integrate with the GCR inference pipeline. Section~\ref{sec:threshold} covers the threshold selection procedure. Section~\ref{sec:scope} defines the scope and boundaries of the implementation. Section~\ref{sec:flowchart} presents the end-to-end system flowchart.

% -----------------------------------------------------------------------------
\section{System Architecture}
\label{sec:architecture}
% -----------------------------------------------------------------------------

The DCA-Trie system is designed as an extension of the Graph-Constrained Reasoning (GCR) framework~\citep{luo2024gcr}, inheriting its four-layer pipeline architecture and modifying the constraint oracle layer to incorporate semantic relevance conditioning. The four layers are: entity linking, constraint oracle, constrained decoding, and inductive reasoning.

The \textbf{entity linking layer} identifies the set of question entities $\mathcal{E}_q$ from the natural language input question $q$ using a named entity recognition tool. These entities serve as the starting points for all KG path search operations. This layer is retained from GCR without modification.

The \textbf{constraint oracle layer} is the primary site of DCA-Trie's contribution. In the baseline GCR system, this layer constructs a KG-Trie by breadth-first search over the KG from $\mathcal{E}_q$ to a maximum depth of $L$ hops, encoding all reachable paths as a prefix tree before generation begins. DCA-Trie replaces this with a semantically conditioned oracle. In DCA-Trie v1, semantic filtering is applied at trie construction time, pruning paths whose cosine similarity to the input question falls below a threshold $\tau$. In DCA-Trie v2, the oracle is additionally updated step-wise during generation, re-scoring the KG neighbourhood of each committed entity against the current reasoning context $(q, y_{<t})$ and admitting only relevant neighbours.

The \textbf{constrained decoding layer} applies the binary logit mask derived from the constraint oracle at every token generation step during beam search. This layer enforces structural faithfulness: only tokens that correspond to valid KG path prefixes in the oracle trie can be generated. This layer is structurally identical in GCR and DCA-Trie; the only change is the oracle it queries.

The \textbf{inductive reasoning layer} uses GPT-4o-mini to synthesise the $K$ reasoning paths produced by constrained beam search into a single final answer. This layer is retained from GCR without modification.

Figure~\ref{fig:architecture} illustrates the full system architecture with the DCA-Trie constraint oracle layer highlighted.

\begin{figure}[H]
    \centering
    \begin{tikzpicture}[
        node distance=0.6cm and 1.2cm,
        box/.style={draw, rectangle, rounded corners=4pt, minimum height=1.0cm,
                minimum width=3.5cm, align=center, thick, font=\small},
        contrib/.style={draw, rectangle, rounded corners=4pt, minimum height=1.0cm,
                minimum width=3.5cm, align=center, thick, font=\small,
                fill=teal!12, draw=teal!60!black},
        arrow/.style={-{Stealth[scale=1.1]}, thick},
        label/.style={font=\footnotesize\itshape, text=gray}
        ]
        \node[box]    (el)   {Entity Linking};
        \node[contrib, below=of el] (oracle) {Constraint Oracle\\(DCA-Trie)};
        \node[box,    below=of oracle] (cd) {Constrained Decoding\\(Beam Search, $K$ beams)};
        \node[box,    below=of cd]    (ir) {Inductive Reasoning\\(GPT-4o-mini)};
        \node[box,    below=of ir]    (ans){Final Answer};

        \draw[arrow] (el)     -- node[label, right] {$\mathcal{E}_q$} (oracle);
        \draw[arrow] (oracle) -- node[label, right] {Mask $\mathbf{m}_t$} (cd);
        \draw[arrow] (cd)     -- node[label, right] {$K$ reasoning paths} (ir);
        \draw[arrow] (ir)     -- (ans);

        \node[label, right=0.3cm of oracle] {\textit{← Primary contribution}};

    \end{tikzpicture}
    \caption{DCA-Trie system architecture. The constraint oracle layer (teal) is the primary contribution; all other layers are inherited from GCR~\citep{luo2024gcr}.}
    \label{fig:architecture}
\end{figure}

% -----------------------------------------------------------------------------
\section{Constraint Oracle Layer}
\label{sec:oracle}
% -----------------------------------------------------------------------------

The constraint oracle layer is responsible for computing, at each generation step $t$, the valid token set $\mathcal{W}_{\text{val}}(t) \subseteq \mathcal{V}$ from which the LLM may select the next token. This section defines the formal problem the oracle is designed to solve, introduces the Semantic Irrelevance Ratio as the primary evaluation metric for oracle quality, and describes the sentence transformer relevance scorer that the oracle uses.

% ---- 3.2.1 ------------------------------------------------------------------
\subsection{Formal Problem Definition}
\label{sec:oracle:problem}

Let $\mathcal{G} = (\mathcal{E}, \mathcal{R}, \mathcal{T})$ be a knowledge graph with entity set $\mathcal{E}$, relation set $\mathcal{R}$, and verified triplet set $\mathcal{T} \subseteq \mathcal{E} \times \mathcal{R} \times \mathcal{E}$. Let $\mathcal{V}$ be the LLM's vocabulary. Given an input question $q$ with associated question entities $\mathcal{E}_q \subseteq \mathcal{E}$ identified by entity linking, the constraint oracle must produce at each step $t$ a valid token set satisfying three conditions simultaneously:

\begin{enumerate}[label=(\roman*)]
    \item \textbf{Structural faithfulness.} Every token $v \in \mathcal{W}_{\text{val}}(t)$ corresponds to a valid prefix of some reasoning path $\pi$ in $\mathcal{G}$. No token outside $\mathcal{G}$ is ever admitted.

    \item \textbf{Semantic relevance.} The proportion of $\mathcal{W}_{\text{val}}(t)$ consisting of semantically irrelevant paths is strictly less than that produced by the baseline GCR oracle.

    \item \textbf{Answer recall.} The gold answer path $\pi^*$ is never excluded from $\mathcal{W}_{\text{val}}(t)$ at any step. The false negative rate on gold answer paths is held below 5\% on the validation set.
\end{enumerate}

Existing frameworks construct the oracle as $\mathcal{W}_{\text{val}}(t) = f(\mathcal{G}, \mathcal{E}_q)$, a static function of graph structure and question entities alone that carries no information about question semantics or the evolving generation context. DCA-Trie upgrades this to:
\[
    \mathcal{W}_{\text{val}}^{\text{DCA}}(t) = f\!\left(\mathcal{G},\, \mathcal{E}_q,\, q,\, y_{<t}\right)
\]
where $q$ is the full input question and $y_{<t}$ is the partial generation up to step $t$. This formulation requires the oracle to maintain and re-evaluate semantic state across generation steps, which is a qualitatively different requirement from static trie traversal.

% ---- 3.2.2 ------------------------------------------------------------------
\subsection{The Semantic Irrelevance Ratio}
\label{sec:oracle:sir}

The Semantic Irrelevance Ratio (SIR) is a new metric introduced in this work to quantify constraint permissiveness independently of downstream answer accuracy. No prior work provides such a metric, making it impossible to directly measure progress on reducing irrelevant paths in the valid set.

% \begin{definition}[Semantic Irrelevance Ratio]
%     Let $\mathcal{T}_{\text{trie}}$ be a KG-Trie and let $\Pi(\mathcal{T}_{\text{trie}})$ be the set of all complete paths encoded in it. Let $\Pi^*(q) \subseteq \Pi(\mathcal{T}_{\text{trie}})$ be the subset of paths that lead to the gold answer entity $a(q)$ for question $q$. The Semantic Irrelevance Ratio is:
%     \[
%         \text{SIR}(q) \;=\; 1 \;-\; \frac{|\Pi^*(q)|}{|\Pi(\mathcal{T}_{\text{trie}})|}
%     \]
% \end{definition}

SIR $\in [0, 1]$. A value of 0 means every path in the trie leads to the gold answer, representing perfect precision. A value close to 1 means almost all paths are irrelevant, placing the full discriminative burden on the LLM's parametric knowledge. For a three-hop CWQ question with Freebase average out-degree $d \approx 20$ and up to three question entities, $|\Pi(\mathcal{T}_{\text{trie}})| \leq 24{,}000$, giving SIR $\approx 0.9999$ under the GCR baseline.

The aggregate SIR over a dataset $\mathcal{D}$ is the macro-average:
\[
    \overline{\text{SIR}} \;=\; \frac{1}{|\mathcal{D}|} \sum_{q \in \mathcal{D}} \text{SIR}(q)
\]

SIR is additionally computed stratified by hop depth, since permissiveness is expected to increase exponentially with $L$. DCA-Trie's primary design objective is to reduce $\overline{\text{SIR}}$ while maintaining answer recall condition (iii) above.

% ---- 3.2.3 ------------------------------------------------------------------
\subsection{Semantic Relevance Scorer}
\label{sec:oracle:scorer}

The semantic relevance scorer is the component that enables the oracle to condition on question semantics and generation context. It is implemented using the \texttt{all-MiniLM-L6-v2} sentence transformer~\citep{reimers2019sbert}, selected for its balance of semantic quality (STS-B score 82.3) and inference speed ($\approx$80ms per query on CPU), which is essential for step-wise application during beam search decoding.

Given a candidate KG path $p = (e_0, r_1, e_1, \ldots, r_L, e_L)$ and the current generation context $(q, y_{<t})$, the relevance score is:
\[
    \text{score}(p,\, q,\, y_{<t}) \;=\; \cos\!\Bigl(E(p),\; E\!\left(\text{concat}(q,\, y_{<t})\right)\Bigr)
\]
where $E(\cdot)$ is the sentence transformer encoder, $\text{concat}(q, y_{<t})$ concatenates the question and partial generation as a single string separated by a delimiter token, and $\cos(\cdot, \cdot)$ is cosine similarity.

\textbf{Path string representation.} A KG path is converted to a natural language string for encoding by replacing Freebase namespace prefixes and underscores with spaces. For example, the path $(\textit{Blue\_Hawaii},\, \textit{film.film.featured\_film\_location},\, \textit{Hawaii})$ becomes the string ``Blue Hawaii featured film location Hawaii''. This representation allows the pretrained encoder to apply its semantic knowledge to assess path relevance without domain-specific fine-tuning.

\textbf{Binary admission criterion.} Paths are admitted to the trie if and only if their relevance score meets or exceeds a threshold $\tau$:
\[
    p \in \mathcal{W}_{\text{val}}^{\text{DCA}}(t) \;\iff\; \text{score}(p,\, q,\, y_{<t}) \;\geq\; \tau
\]
This converts the continuous cosine similarity signal into a hard binary constraint, consistent with the logit-masking paradigm: paths below $\tau$ have their corresponding token logits set to $-\infty$.

% -----------------------------------------------------------------------------
\section{DCA-Trie v1: Static Semantic Filtering}
\label{sec:v1}
% -----------------------------------------------------------------------------

DCA-Trie v1 is the static variant of the framework. It applies the relevance scorer once at trie construction time, filtering paths by their similarity to the input question $q$ before generation begins. This preserves the computational efficiency of GCR's precomputed oracle while reducing SIR by removing question-irrelevant paths from the trie.

% ---- 3.3.1 ------------------------------------------------------------------
\subsection{Construction Procedure}

DCA-Trie v1 modifies GCR's KG-Trie construction by inserting a semantic filtering step between BFS path enumeration and trie insertion. The procedure is formalised in Algorithm~\ref{alg:v1}.

\begin{figure}[H]
    \begin{algorithm}[H]
        \caption{DCA-Trie v1 --- Static Semantic Filtering}
        \label{alg:v1}
        \begin{algorithmic}[1]
            \Require Question $q$, entity set $\mathcal{E}_q$, KG $\mathcal{G}$, hop depth $L$, threshold $\tau$
            \Ensure Filtered trie $\mathcal{T}_{\text{filtered}}$
            \State $\Pi_{\text{all}} \leftarrow \textsc{BFS}(\mathcal{G},\, \mathcal{E}_q,\, L)$
            \Comment{Enumerate all $L$-hop paths, as in GCR}
            \State $\mathbf{c} \leftarrow E(q)$
            \Comment{Encode question once}
            \State $\Pi_{\text{filtered}} \leftarrow \emptyset$
            \For{each path $\pi \in \Pi_{\text{all}}$}
            \State $s \leftarrow \cos\!\bigl(E(\text{path\_string}(\pi)),\; \mathbf{c}\bigr)$
            \If{$s \geq \tau$}
            \State $\Pi_{\text{filtered}} \leftarrow \Pi_{\text{filtered}} \cup \{\pi\}$
            \EndIf
            \EndFor
            \State $\mathcal{T}_{\text{filtered}} \leftarrow \textsc{BuildTrie}(\Pi_{\text{filtered}})$
            \State \Return $\mathcal{T}_{\text{filtered}}$
        \end{algorithmic}
    \end{algorithm}
    \caption{Algorithm for DCA-Trie v1 static semantic filtering at trie construction time.}
    \label{fig:alg_v1}
\end{figure}

% ---- 3.3.2 ------------------------------------------------------------------
\subsection{Properties and Limitations}

DCA-Trie v1 satisfies all three conditions of the problem definition in Section~\ref{sec:oracle:problem}. Structural faithfulness is preserved because the filter can only remove paths from GCR's valid set, never introduce new ones; every path in $\mathcal{T}_{\text{filtered}}$ is a verified KG path. Semantic relevance improves because $|\Pi_{\text{filtered}}| \leq |\Pi_{\text{all}}|$, reducing SIR when the threshold $\tau$ is set appropriately. Answer recall is preserved to the extent that gold answer paths score above $\tau$, which is controlled by the threshold selection procedure in Section~\ref{sec:threshold}.

Because $\mathcal{T}_{\text{filtered}}$ is precomputed before generation begins, DCA-Trie v1 preserves the amortised constant-time trie query efficiency established by~\citet{willard2023outlines}. The oracle is queried identically to GCR at inference time; the only difference is the size of the trie being queried.

The limitation of DCA-Trie v1 is that filtering is conditioned on $q$ alone. The trie is static once constructed, identical at every generation step. It addresses the question-semantic dimension of context conditioning but not the generation-state dimension: it cannot narrow the valid set further as reasoning commits to a specific direction during generation.

% -----------------------------------------------------------------------------
\section{DCA-Trie v2: Step-Wise Dynamic Expansion}
\label{sec:v2}
% -----------------------------------------------------------------------------

DCA-Trie v2 implements the full context-conditioned oracle formulation $f(\mathcal{G}, \mathcal{E}_q, q, y_{<t})$. It extends DCA-Trie v1 by additionally updating the oracle at each generation step when an entity is committed, re-scoring the KG neighbourhood of the committed entity against the current context $(q, y_{<t})$ and admitting only semantically relevant neighbours.

% ---- 3.4.1 ------------------------------------------------------------------
\subsection{Step-Wise Update Procedure}

DCA-Trie v2 begins with the trie produced by DCA-Trie v1's static filtering as its initial oracle. Each time the LLM commits to an entity token $e^*$ during beam search, the trie is updated by scoring the KG neighbourhood $\mathcal{N}(e^*)$ against the updated context and admitting qualifying neighbours. Algorithm~\ref{alg:v2} formalises this procedure.

\begin{figure}[H]
    \begin{algorithm}[H]
        \caption{DCA-Trie v2 --- Step-Wise Dynamic Expansion}
        \label{alg:v2}
        \begin{algorithmic}[1]
            \Require Question $q$, entity set $\mathcal{E}_q$, KG $\mathcal{G}$, hop depth $L$, threshold $\tau$
            \Ensure Updated tries $\mathcal{T}_0, \mathcal{T}_1, \ldots, \mathcal{T}_T$
            \State $\mathcal{T}_0 \leftarrow \textsc{DCA-Trie-v1}(q,\, \mathcal{E}_q,\, \mathcal{G},\, L,\, \tau)$
            \For{each generation step $t = 1, 2, \ldots, T$}
            \State Apply mask from $\mathcal{T}_{t-1}$ during beam search to select token $y_t$
            \If{$y_t$ is an entity token $e^*$}
            \State $\text{context} \leftarrow \text{concat}(q,\; y_{\leq t})$
            \State $\mathbf{c}_t \leftarrow E(\text{context})$
            \State $\mathcal{N}(e^*) \leftarrow \{(e^*, r, e') \in \mathcal{G} \mid \text{for all } r, e'\}$
            \State $\mathcal{N}_{\text{admitted}} \leftarrow \emptyset$
            \For{each triplet $(e^*, r, e') \in \mathcal{N}(e^*)$}
            \State $s \leftarrow \cos\!\bigl(E(\text{path\_string}(e^*, r, e')),\; \mathbf{c}_t\bigr)$
            \If{$s \geq \tau$}
            \State $\mathcal{N}_{\text{admitted}} \leftarrow \mathcal{N}_{\text{admitted}} \cup \{(e^*, r, e')\}$
            \EndIf
            \EndFor
            \State $\mathcal{T}_t \leftarrow \mathcal{T}_{t-1} \;\cup\; \textsc{BuildTrie}(\mathcal{N}_{\text{admitted}})$
            \Else
            \State $\mathcal{T}_t \leftarrow \mathcal{T}_{t-1}$
            \EndIf
            \EndFor
        \end{algorithmic}
    \end{algorithm}
    \caption{Algorithm for DCA-Trie v2 step-wise dynamic expansion with semantic filtering.}
    \label{fig:alg_v2}
\end{figure}

% ---- 3.4.2 ------------------------------------------------------------------
\subsection{Comparison with DoG}

DCA-Trie v2's step-wise expansion is directly comparable to Decoding on Graphs (DoG;~\citealt{li2024dog}). The structural difference is the admission condition:

\begin{align*}
    \text{DoG:}         & \quad \mathcal{T}_{t+1} = \mathcal{T}_t \cup \{(e^*, r, e') : (e^*, r, e') \in \mathcal{G}\}                                                                  \\
    \text{DCA-Trie v2:} & \quad \mathcal{T}_{t+1} = \mathcal{T}_t \cup \{(e^*, r, e') : (e^*, r, e') \in \mathcal{G} \;\text{and}\; \text{score}(e^*re',\, q,\, y_{\leq t}) \geq \tau\}
\end{align*}

DoG admits every KG neighbour of each committed entity unconditionally. DCA-Trie v2 adds a single semantic scoring condition. This change is the core algorithmic contribution of the step-wise variant: the structural mechanism of step-wise expansion is inherited from DoG, but the semantic filter converts topologically dynamic but semantically blind expansion into context-conditioned expansion.

% ---- 3.4.3 ------------------------------------------------------------------
\subsection{Efficiency Analysis}

DCA-Trie v2 incurs two costs beyond GCR that do not appear in DCA-Trie v1: step-wise trie updates and per-step relevance scoring. The trie update cost at each entity step is $O(|\mathcal{N}_{\text{admitted}}| \times L)$, where $|\mathcal{N}_{\text{admitted}}| \ll |\mathcal{N}(e^*)|$ when $\tau$ is well-calibrated, substantially reducing the expansion cost relative to DoG. The scoring cost is one encoder forward pass per candidate neighbour path, approximately 80ms per query on CPU, amortised across the beam width $K$ because the context embedding $\mathbf{c}_t$ is computed once per step and reused. Over a full generation sequence of $L$ entity steps with beam width $K$, the additional cost relative to GCR is bounded and measurable in the evaluation experiments.

% -----------------------------------------------------------------------------
\section{Integration with the GCR Inference Pipeline}
\label{sec:integration}
% -----------------------------------------------------------------------------

Both DCA-Trie variants are implemented as extensions to GCR's publicly released inference pipeline~\citep{luo2024gcr}. This section describes how each variant integrates with GCR's existing components without modifying the LLM backbone, beam search procedure, or inductive reasoning layer.

% ---- 3.5.1 ------------------------------------------------------------------
\subsection{Pipeline Integration}

GCR's inference pipeline consists of four stages: entity linking using the ELQ entity linker~\citep{decao2021autoregressive} to identify $\mathcal{E}_q$; KG-Trie construction by BFS over preprocessed Freebase subgraphs; constrained beam search using GCR's released Llama-3.1-8B checkpoint~\citep{dubey2024llama3} with the trie queried at each token step; and inductive reasoning using GPT-4o-mini to synthesise $K=8$ beam hypotheses.

DCA-Trie v1 modifies only stage 2: BFS path enumeration is followed by Algorithm~\ref{alg:v1}'s semantic filtering, producing $\mathcal{T}_{\text{filtered}}$ in place of $\mathcal{T}_{\text{GCR}}$. All other stages are unchanged. DCA-Trie v2 additionally modifies stage 3: the beam search loop is augmented with an entity-step hook that triggers Algorithm~\ref{alg:v2}'s trie update whenever an entity token is committed. No modification is made to the LLM backbone, the beam search scoring function, or the inductive reasoning component.

% ---- 3.5.2 ------------------------------------------------------------------
\subsection{Faithfulness Verification}

Structural faithfulness is verified after generation by exhaustively checking every triplet in each generated reasoning chain against GCR's preprocessed Freebase subgraphs. A generated chain is structurally faithful if and only if every triplet $(e_{i-1}, r_i, e_i)$ in the chain exists in the trie oracle used for that generation. This verification procedure follows GCR's original faithfulness evaluation protocol and produces the faithfulness rate reported alongside accuracy results.

% -----------------------------------------------------------------------------
\section{Threshold Selection}
\label{sec:threshold}
% -----------------------------------------------------------------------------

The relevance threshold $\tau$ controls the precision-recall tradeoff of the semantic filter and is the primary hyperparameter of both DCA-Trie variants.

% ---- 3.6.1 ------------------------------------------------------------------
\subsection{Selection Procedure}

The threshold is selected on a held-out validation subset of 100 questions from WebQSP, ensuring no test-set leakage. For each candidate $\tau \in \{0.1, 0.2, 0.3, 0.4, 0.5, 0.6, 0.7, 0.8\}$, two quantities are measured:

\begin{enumerate}[label=(\roman*)]
    \item \textbf{SIR reduction}: the percentage decrease in average SIR relative to the GCR baseline.
    \item \textbf{False negative rate (FNR)}: the proportion of questions for which the gold answer path is excluded from $\mathcal{T}_{\text{filtered}}$ at any step.
\end{enumerate}

The selected threshold is the highest $\tau$ that maintains FNR $< 5\%$ on the validation subset. This criterion prioritises answer recall over SIR reduction: excluding a correct answer path is a recall failure that cannot be recovered downstream, while admitting irrelevant paths is a precision failure that the LLM's residual parametric knowledge can compensate for. The same selected $\tau$ is applied to both DCA-Trie v1 and v2, enabling a controlled comparison.

% ---- 3.6.2 ------------------------------------------------------------------
\subsection{Ablation over Threshold Values}

In addition to the selected threshold, an ablation study evaluates DCA-Trie's sensitivity to $\tau$ by reporting full evaluation metrics for $\tau \in \{0.3, 0.4, 0.5, 0.6\}$ on the WebQSP test set. This directly addresses the question of whether the SIR-accuracy tradeoff is stable across a range of thresholds and whether the selected value is a local or global optimum on the precision-recall curve.

% -----------------------------------------------------------------------------
\section{Scope of Work}
\label{sec:scope}
% -----------------------------------------------------------------------------

The following components are within the implementation scope of this project. DCA-Trie v1 and v2 are implemented and evaluated on WebQSP and CWQ benchmarks against GCR and chain-of-thought baselines. The SIR metric is defined, computed, and reported stratified by hop depth. The threshold selection procedure is implemented on a held-out validation subset. An ablation study over $\tau$ values is conducted on the WebQSP test set.

The following are explicitly outside scope and identified as directions for future work. Confidence-conditioned dynamic threshold adjustment, which would modulate $\tau$ based on LLM token entropy at each step, is a natural extension not implemented in this work. Fine-tuning the sentence transformer scorer on KGQA-specific sentence pairs would be expected to improve filtering precision but requires labelled training data beyond what is used here. Generalisation of DCA-Trie to knowledge graphs other than Freebase --- such as Wikidata or domain-specific KGs --- is architecturally straightforward but not evaluated. Production-scale deployment and inference latency optimisation are engineering extensions beyond the research scope.

All code, experimental configurations, and evaluation data will be released publicly to support reproducibility of all reported findings.

% -----------------------------------------------------------------------------
\section{End-to-End System Flowchart}
\label{sec:flowchart}
% -----------------------------------------------------------------------------

Figure~\ref{fig:flowchart} illustrates the complete end-to-end processing flow from input question to final answer, showing the decision points specific to DCA-Trie v1 and v2.

\begin{figure}[H]
    \centering
    \begin{tikzpicture}[
        node distance=0.55cm,
        startstop/.style={draw, rectangle, minimum height=0.85cm,
                minimum width=3.2cm, align=center, thick, font=\small},
        process/.style={draw, rectangle, minimum height=0.85cm,
                minimum width=3.8cm, align=center, thick, font=\small},
        contrib/.style={draw, rectangle, minimum height=0.85cm,
                minimum width=3.8cm, align=center, thick, font=\small,
                fill=teal!12, draw=teal!60!black},
        decision/.style={draw, diamond, minimum height=1.0cm, minimum width=1.0cm,
                align=center, thick, font=\scriptsize, aspect=2.5},
        arrow/.style={-{Stealth[scale=1.0]}, thick},
        label/.style={font=\scriptsize, midway}
        ]
        \node[startstop]  (start)  {Input question $q$};
        \node[process,  below=of start]  (el)   {Entity linking\\identify $\mathcal{E}_q$};
        \node[process,  below=of el]     (bfs)  {BFS over $\mathcal{G}$ from $\mathcal{E}_q$\\enumerate $\Pi_{\text{all}}$};
        \node[contrib,  below=of bfs]    (sf)   {Semantic filtering (v1)\\score paths vs.\ $q$; keep $\geq \tau$};
        \node[process,  below=of sf]     (trie) {Build filtered trie $\mathcal{T}_{\text{filtered}}$};
        \node[process,  below=of trie]   (beam) {Beam search step $t$:\\apply mask from $\mathcal{T}_{t-1}$, select $y_t$};
        \node[decision, below=of beam]   (ent)  {$y_t$ is\\entity?};
        \node[contrib,  right=2.5cm of ent] (update) {Re-score $\mathcal{N}(e^*)$\\vs.\ $(q, y_{\leq t})$; update $\mathcal{T}_t$ (v2)};
        \node[decision, below=of ent]   (done) {Generation\\complete?};
        \node[process,  below=of done]   (ir)   {Inductive reasoning\\synthesise $K$ paths};
        \node[startstop, below=of ir]   (ans)  {Final answer};

        \draw[arrow] (start) -- (el);
        \draw[arrow] (el)    -- (bfs);
        \draw[arrow] (bfs)   -- (sf);
        \draw[arrow] (sf)    -- (trie);
        \draw[arrow] (trie)  -- (beam);
        \draw[arrow] (beam)  -- (ent);
        \draw[arrow] (ent)   -- node[label, right] {Yes} (update);
        \draw[arrow] (update) |- (done);
        \draw[arrow] (ent)   -- node[label, left]  {No}  (done);
        \draw[arrow] (done)  -- node[label, right] {Yes} (ir);
        \draw[arrow] (done.west) -- ++(-1.5,0) |- node[label, left, pos=0.25] {No} (beam.west);
        \draw[arrow] (ir)    -- (ans);

    \end{tikzpicture}
    \caption{End-to-end DCA-Trie processing flowchart. Teal boxes indicate the primary contributions of this work. The v2 trie update step is conditional on an entity token being committed at step $t$.}
    \label{fig:flowchart}
\end{figure}
